% ============================================================
%  C++ and MLIR Fast Track
%  Compile: cd study && xelatex book.tex && xelatex book.tex
% ============================================================
\documentclass[11pt,a4paper]{book}

% --- Fonts & Korean ---
\usepackage{fontspec}
\usepackage{kotex}
\setmainfont{Palatino}
\setsansfont{Apple SD Gothic Neo}
\setmonofont{Menlo}[Scale=0.85]
\setmainhangulfont{Apple SD Gothic Neo}
\setsanshangulfont{Apple SD Gothic Neo}
\setmonohangulfont{AppleGothic}[Scale=0.92]

% --- Layout ---
\usepackage[margin=2.5cm]{geometry}
\usepackage{setspace}
\onehalfspacing

% --- Packages ---
\usepackage{listings}
\usepackage{xcolor}
\usepackage{tcolorbox}
\tcbuselibrary{skins,breakable}
\usepackage{enumitem}
\usepackage{hyperref}
\usepackage{titlesec}
\usepackage{fancyhdr}
\usepackage{graphicx}
\usepackage{booktabs}
\usepackage{longtable}
\usepackage{array}
\usepackage{tabularx}

% --- Colors ---
\definecolor{codebg}{HTML}{F7F6F2}
\definecolor{codeframe}{HTML}{C8C8C8}
\definecolor{keyword}{HTML}{0B57D0}
\definecolor{string}{HTML}{A31515}
\definecolor{comment}{HTML}{2E7D32}
\definecolor{type}{HTML}{1565C0}
\definecolor{boxblue}{HTML}{1565C0}
\definecolor{boxgreen}{HTML}{2E7D32}
\definecolor{boxorange}{HTML}{E65100}
\definecolor{boxpurple}{HTML}{6A1B9A}
\definecolor{quizbg}{HTML}{FFF8E1}
\definecolor{summarybg}{HTML}{E3F2FD}
\definecolor{checkbg}{HTML}{E8F5E9}

% --- Code listings ---
\lstset{
  basicstyle=\small\ttfamily,
  backgroundcolor=\color{codebg},
  frame=single,
  rulecolor=\color{codeframe},
  breaklines=true,
  breakatwhitespace=false,
  tabsize=2,
  showstringspaces=false,
  numbers=left,
  numberstyle=\tiny\color{gray},
  numbersep=8pt,
  aboveskip=10pt,
  belowskip=10pt,
  xleftmargin=15pt,
  framexleftmargin=15pt,
  columns=flexible,
  keepspaces=true,
}

\lstdefinestyle{cmake}{
  language={},
  keywordstyle=\color{keyword}\bfseries,
  morekeywords={
    cmake_minimum_required,project,add_executable,add_library,target_link_libraries,
    target_include_directories,target_compile_features,target_compile_definitions,
    target_compile_options,target_sources,target_link_options,target_link_directories,
    find_package,include,if,else,endif,set,option,message,install,export,include_guard,
    add_subdirectory,enable_testing,add_test,set_property,set_target_properties,
    configure_file,file,list,foreach,endforeach,function,endfunction,macro,endmacro,
    PUBLIC,PRIVATE,INTERFACE,REQUIRED,CONFIG,VERSION,LANGUAGES
  },
  stringstyle=\color{string},
  commentstyle=\color{comment}\itshape,
  morecomment=[l]{\#},
}

\lstdefinestyle{cpp}{
  language=C++,
  keywordstyle=\color{keyword}\bfseries,
  stringstyle=\color{string},
  commentstyle=\color{comment}\itshape,
  morekeywords={override,final,nullptr,constexpr,noexcept,auto,decltype,static_assert,
                unique_ptr,shared_ptr,make_unique,make_shared,optional,string_view},
  emph={std,string,vector,map,cout,cerr,endl,optional,unique_ptr},
  emphstyle=\color{type},
}

\lstdefinestyle{shell}{
  language=bash,
  keywordstyle=\color{keyword}\bfseries,
  stringstyle=\color{string},
  commentstyle=\color{comment}\itshape,
}

% --- Boxes ---
\newtcolorbox{summarybox}[1][]{
  colback=summarybg,
  colframe=boxblue,
  fonttitle=\bfseries,
  title={배운 것 정리},
  breakable,
  #1
}

\newtcolorbox{quizbox}[1][]{
  colback=quizbg,
  colframe=boxorange,
  fonttitle=\bfseries,
  title={Quiz},
  breakable,
  #1
}

\newtcolorbox{conceptbox}[1][]{
  colback=white,
  colframe=boxgreen,
  fonttitle=\bfseries,
  breakable,
  #1
}

\newtcolorbox{checklistbox}[1][]{
  colback=checkbg,
  colframe=boxgreen,
  fonttitle=\bfseries,
  title={실전 체크리스트},
  breakable,
  #1
}

\newtcolorbox{warningbox}[1][]{
  colback=white,
  colframe=boxorange,
  fonttitle=\bfseries,
  title={주의},
  breakable,
  #1
}

\newtcolorbox{pythonbox}[1][]{
  colback=white,
  colframe=boxpurple,
  fonttitle=\bfseries,
  title={Python 개발자 관점},
  breakable,
  #1
}

% --- Header/Footer ---
\pagestyle{fancy}
\fancyhf{}
\fancyhead[LE]{\leftmark}
\fancyhead[RO]{\rightmark}
\fancyfoot[C]{\thepage}
\renewcommand{\headrulewidth}{0.4pt}
\setlength{\headheight}{14pt}

% --- Chapter title format ---
\titleformat{\chapter}[display]
  {\normalfont\huge\bfseries}
  {\chaptertitlename\ \thechapter}{20pt}{\Huge}

% --- Hyperref ---
\hypersetup{
  colorlinks=true,
  linkcolor=boxblue,
  urlcolor=boxblue,
  pdftitle={C++ and MLIR Fast Track},
  pdfauthor={Codex},
}

% --- Convenience macros ---
\DeclareRobustCommand{\code}[1]{\texttt{\detokenize{#1}}}
\DeclareRobustCommand{\file}[1]{\texttt{\detokenize{#1}}}
\newcommand{\term}[1]{\textbf{#1}}

\begin{document}

\begin{titlepage}
\centering
\vspace*{2.0cm}
{\Huge\bfseries C++ and MLIR Fast Track\par}
\vspace{0.6cm}
{\LARGE C++ 기본 개념, 실전 패턴, MLIR 감각까지 빠르게 익히는 교재\par}
\vspace{1.4cm}
{\Large\itshape Python 개발자가 C++/MLIR에 빠르게 익숙해지기 위한 압축 학습서\par}
\vspace{2.2cm}
{\large 목표\par}
\vspace{0.3cm}
{\large 1. CMake와 빌드 흐름을 읽을 수 있다.\par}
{\large 2. C++ 기본 개념을 설계 판단 수준으로 이해한다.\par}
{\large 3. 자주 쓰는 코드 패턴을 통째로 외워 실전에 바로 적용한다.\par}
{\large 4. MLIR 프로젝트 구조와 개발 흐름을 큰 그림으로 이해한다.\par}
\vfill
{\large 2026\par}
\end{titlepage}

\frontmatter
\tableofcontents

\mainmatter

\chapter{이 책을 읽는 방법}

이 책의 목적은 C++ 전체를 완벽히 가르치는 것이 아니다. 대신 \term{C++ 코드가 왜 저렇게 생겼는지 이해하고}, \term{자주 나오는 패턴을 암기하고}, \term{MLIR 코드를 읽고 수정할 수 있는 상태}까지 빠르게 끌어올리는 것이다.

\begin{conceptbox}[title={권장 학습 순서}]
\begin{enumerate}[nosep]
\item Part I에서 CMake와 빌드 흐름을 먼저 익힌다.
\item Part II에서 소유권, 템플릿, virtual, move 같은 기본 개념을 잡는다.
\item Part III에서 실전 패턴을 암기한다.
\item Part IV에서 MLIR의 큰 그림, TableGen, 패스 개발 흐름을 연결한다.
\end{enumerate}
\end{conceptbox}

\begin{pythonbox}
Python에서는 런타임 모델이 대부분을 숨겨준다. C++에서는 타입, 수명, 소유권, 링크, 빌드 그래프를 직접 다뤄야 한다. 그래서 C++를 잘하려면 문법보다 먼저 \term{객체 수명}과 \term{타겟 관계}를 이해해야 한다.
\end{pythonbox}

\part{빌드와 CMake}

\chapter{MLIR 프로젝트를 위한 CMake 핵심}

LLVM/MLIR을 다루려면 먼저 \code{cmake}가 무엇을 하는지 감이 잡혀야 한다. 핵심은 \term{소스 파일 나열}이 아니라 \term{타겟 그래프 설계}다.

\section{전체 흐름}

\begin{lstlisting}[style=shell,numbers=none,title={LLVM/MLIR 빌드의 기본 흐름}]
git clone https://github.com/llvm/llvm-project.git
cd llvm-project
git checkout llvmorg-18.1.0

cmake -S llvm -B build \
  -DCMAKE_BUILD_TYPE=Release \
  -DCMAKE_INSTALL_PREFIX=$HOME/llvm-install \
  -DLLVM_ENABLE_PROJECTS=mlir \
  -DLLVM_TARGETS_TO_BUILD=host \
  -DLLVM_ENABLE_ASSERTIONS=ON

cmake --build build -j$(sysctl -n hw.ncpu)
cmake --install build
\end{lstlisting}

\begin{summarybox}
\begin{itemize}[nosep]
\item \code{-S}는 소스 루트, \code{-B}는 빌드 디렉토리다.
\item \code{configure} 단계와 \code{build} 단계는 다르다.
\item MLIR 개발에서는 \code{Release + Assertions} 조합이 실용적이다.
\item LLVM은 변동이 빠르므로 \code{main}보다 릴리스 태그 고정이 안전하다.
\end{itemize}
\end{summarybox}

\section{외워야 할 CMake 명령 5개}

\begin{lstlisting}[style=cmake,title={가장 자주 보는 CMake 명령}]
cmake_minimum_required(VERSION 3.20)
project(MyTool LANGUAGES CXX)

add_library(MyDialect lib/MyDialect.cpp)
add_executable(my-opt tools/my-opt.cpp)

target_link_libraries(my-opt PRIVATE MyDialect MLIRIR)

find_package(MLIR REQUIRED CONFIG)
\end{lstlisting}

\begin{conceptbox}[title={각 명령의 역할}]
\begin{itemize}[nosep]
\item \code{project}: 프로젝트 이름과 언어를 선언한다.
\item \code{add_library}: 재사용 가능한 타겟을 만든다.
\item \code{add_executable}: 최종 바이너리를 만든다.
\item \code{target_link_libraries}: 링크 의존성과 사용 전파를 선언한다.
\item \code{find_package}: 외부 패키지의 경로와 설정을 불러온다.
\end{itemize}
\end{conceptbox}

\section{\code{find_package}, \code{include}, \code{PUBLIC/PRIVATE}}

\begin{lstlisting}[style=cmake,title={MLIR 프로젝트 최소 골격}]
find_package(LLVM REQUIRED CONFIG)
find_package(MLIR REQUIRED CONFIG)

list(APPEND CMAKE_MODULE_PATH "${LLVM_CMAKE_DIR}")
list(APPEND CMAKE_MODULE_PATH "${MLIR_CMAKE_DIR}")

include(AddLLVM)
include(AddMLIR)

add_executable(my-opt tools/my-opt.cpp)
target_link_libraries(my-opt PRIVATE MLIRIR MLIRParser MLIRPass)
\end{lstlisting}

\begin{pythonbox}
\code{find_package(MLIR)}는 Python의 \code{import mlir}와 완전히 같지는 않지만, 감각적으로는 비슷하다. MLIR이 어디 설치되어 있고 어떤 include/lib/cmake 파일을 써야 하는지 현재 프로젝트로 가져오는 단계다.
\end{pythonbox}

\begin{warningbox}
\code{PUBLIC}과 \code{PRIVATE}를 대충 쓰면 의존성이 엉켜서 나중에 디버깅이 어렵다. 라이브러리의 공개 API에 필요한 것은 \code{PUBLIC}, 내부 구현에서만 쓰는 것은 \code{PRIVATE}로 본다.
\end{warningbox}

\section{실전에서 자주 나는 에러}

\begin{longtable}{p{0.28\textwidth}p{0.28\textwidth}p{0.34\textwidth}}
\toprule
증상 & 원인 후보 & 먼저 볼 곳 \\
\midrule
\endhead
\code{undefined reference} & 링크 라이브러리 누락 & \code{target_link_libraries} \\
\code{file not found} & include 경로 누락 & \code{target_include_directories}, \code{find_package} \\
\code{No rule to make target} & 타겟/생성 파일 연결 누락 & \code{add_subdirectory}, generated sources \\
\code{RTTI mismatch} & 컴파일 옵션 불일치 & LLVM/MLIR 빌드 옵션, \code{-fno-rtti} 여부 \\
\bottomrule
\end{longtable}

\begin{checklistbox}
\begin{itemize}[nosep]
\item 소스와 빌드를 분리하는 out-of-source build를 유지할 것
\item 실행파일보다 라이브러리 단위로 구조를 나눌 것
\item MLIR을 찾지 못하면 먼저 \code{MLIR_DIR}과 \code{LLVM_DIR}를 의심할 것
\item 에러가 링크 단계인지 컴파일 단계인지 먼저 분리할 것
\end{itemize}
\end{checklistbox}

\part{C++ 기본 개념}

\chapter{C++를 읽을 때 가장 먼저 봐야 하는 것}

C++ 코드를 읽을 때 변수 타입보다 먼저 봐야 하는 것은 \term{누가 소유하는가}, \term{언제 파괴되는가}, \term{런타임 디스패치인가 컴파일 타임 디스패치인가}다.

\section{Template vs Virtual}

\begin{conceptbox}[title={한 줄 요약}]
\begin{itemize}[nosep]
\item \term{Virtual}: 런타임에 어떤 자식이 올지 모를 때
\item \term{Template}: 컴파일 타임에 타입이 이미 정해져 있을 때
\end{itemize}
\end{conceptbox}

\begin{lstlisting}[style=cpp,title={Virtual: 런타임 디스패치}]
class Serializer {
public:
    virtual ~Serializer() = default;
    virtual std::string name() const = 0;
};

class JsonSerializer : public Serializer {
public:
    std::string name() const override { return "json"; }
};

void use(Serializer &s) {
    std::cout << s.name() << "\n";
}
\end{lstlisting}

\begin{lstlisting}[style=cpp,title={Template: 컴파일 타임 디스패치}]
template <typename T>
void dump(const T &value) {
    std::cout << value.toString() << "\n";
}
\end{lstlisting}

\begin{summarybox}
\begin{itemize}[nosep]
\item factory가 부모 포인터를 반환하면 virtual일 가능성이 높다.
\item \code{OpConversionPattern<OpT>}처럼 타입이 템플릿 인자로 박혀 있으면 template 쪽이다.
\item 성능보다 설계 이유를 먼저 본다. 런타임 선택이 필요하면 virtual이 맞다.
\end{itemize}
\end{summarybox}

\section{포인터와 참조: 소유권 관점으로 보기}

\begin{longtable}{p{0.22\textwidth}p{0.2\textwidth}p{0.48\textwidth}}
\toprule
형태 & 소유권 & 언제 쓰나 \\
\midrule
\endhead
\code{T*} & 없음 & null 가능, 빌려서 잠깐 접근 \\
\code{T\&} & 없음 & null이 되면 안 되는 별명 \\
\code{std::unique_ptr<T>} & 단독 소유 & 수명 관리가 필요한 힙 객체 \\
\code{std::shared_ptr<T>} & 공유 소유 & 정말 여러 주체가 함께 오래 살아야 할 때 \\
\bottomrule
\end{longtable}

\begin{pythonbox}
Python에서는 이름이 객체를 참조하고 GC가 나중에 정리한다. C++에서는 \term{누가 해제 책임을 갖는지}를 코드 구조에 드러내야 한다. 그래서 \code{unique_ptr}이 자주 나오고, raw pointer는 보통 ``내 것이 아니다''라는 신호다.
\end{pythonbox}

\section{Move Semantics}

\begin{lstlisting}[style=cpp,title={std::move의 의미}]
std::string name = "alice";
std::vector<std::string> names;

names.push_back(name);            // 복사
names.push_back(std::move(name)); // 이동
\end{lstlisting}

\begin{conceptbox}[title={중요한 오해 바로잡기}]
\code{std::move}는 실제로 이동하지 않는다. ``이 객체를 rvalue로 취급해도 된다''고 표시하는 캐스트에 가깝다. 실제 이동은 이동 생성자나 이동 대입 연산자가 수행한다.
\end{conceptbox}

\section{Value Semantics vs Reference Semantics}

\begin{lstlisting}[style=cpp,title={복사를 피하는 기본 습관}]
void process(const std::vector<int> &data); // 읽기 전용, 복사 없음
void mutate(std::vector<int> &data);        // 수정 가능
void consume(std::vector<int> data);        // 값으로 받으면 소유권 복사
\end{lstlisting}

\begin{summarybox}
\begin{itemize}[nosep]
\item 큰 객체를 읽기만 하면 \code{const T\&}
\item 수정해야 하면 \code{T\&}
\item 소유권을 넘기거나 복사본이 필요하면 값 전달
\item 작은 값 타입은 그냥 값으로 받아도 된다
\end{itemize}
\end{summarybox}

\section{LLVM 스타일 타입 시스템}

LLVM/MLIR 코드에서는 C++ RTTI 대신 \code{isa}, \code{cast}, \code{dyn\_cast}를 많이 쓴다.

\begin{lstlisting}[style=cpp,title={LLVM 캐스팅 3형제}]
if (isa<mlir::FuncOp>(op)) {
    auto func = cast<mlir::FuncOp>(op);
}

if (auto func = dyn_cast<mlir::FuncOp>(op)) {
    // 성공했을 때만 func 사용
}
\end{lstlisting}

\begin{warningbox}
\code{cast<T>}는 실패하면 죽는다. 성공이 확실할 때만 쓴다. 확실하지 않으면 \code{dyn\_cast<T>}부터 쓴다.
\end{warningbox}

\chapter{Python 개발자가 헷갈리기 쉬운 C++ 지점}

\section{헤더와 구현 분리}

\begin{lstlisting}[style=cpp,title={헤더 파일}]
#ifndef LOGGER_H
#define LOGGER_H

#include <string>

class Logger {
public:
    void info(const std::string &msg);
};

#endif
\end{lstlisting}

\begin{lstlisting}[style=cpp,title={구현 파일}]
#include "Logger.h"
#include <iostream>

void Logger::info(const std::string &msg) {
    std::cout << msg << "\n";
}
\end{lstlisting}

\begin{pythonbox}
Python은 파일 자체가 모듈이다. C++은 선언과 정의를 분리하면서 컴파일 시간을 줄이고, 여러 번 포함되는 상황을 제어한다.
\end{pythonbox}

\section{Forward Declaration}

\begin{lstlisting}[style=cpp,title={가능한 경우 전방 선언}]
class Engine;

class Car {
public:
    void setEngine(Engine *engine);
private:
    Engine *engine_ = nullptr;
};
\end{lstlisting}

\begin{summarybox}
\begin{itemize}[nosep]
\item 포인터/참조만 들고 있으면 전방 선언으로 충분한 경우가 많다.
\item 객체 자체를 멤버로 들고 있으면 완전한 정의가 필요하다.
\item 헤더 의존성을 줄이면 전체 빌드 시간이 줄어든다.
\end{itemize}
\end{summarybox}

\section{CLI 도구의 기본 골격}

\begin{lstlisting}[style=cpp,title={LLVM 스타일 CLI 도구}]
#include "llvm/Support/CommandLine.h"
#include "llvm/Support/InitLLVM.h"

llvm::cl::opt<std::string> inputFile(
    llvm::cl::Positional,
    llvm::cl::desc("<input file>"),
    llvm::cl::Required);

int main(int argc, char **argv) {
    llvm::InitLLVM y(argc, argv);
    llvm::cl::ParseCommandLineOptions(argc, argv);
    return 0;
}
\end{lstlisting}

\part{암기할 가치가 있는 C++ 패턴}

\chapter{매일 보게 되는 패턴}

\section{RAII}

\begin{lstlisting}[style=cpp,title={자원 획득과 해제를 객체 수명에 묶기}]
{
    std::ifstream file("data.txt");
    std::string line;
    while (std::getline(file, line)) {
        // 사용
    }
} // file 소멸자에서 자동 close
\end{lstlisting}

\begin{conceptbox}[title={핵심}]
RAII는 ``자원을 얻자마자 관리 객체 안에 넣고, 해제는 소멸자에 맡긴다''는 습관이다. C++에서 \code{finally} 대신 소멸자가 중요한 이유다.
\end{conceptbox}

\section{Rule of Zero}

\begin{lstlisting}[style=cpp,title={가능하면 직접 소멸자를 쓰지 않는다}]
struct Config {
    std::string name;
    std::vector<int> values;
    std::unique_ptr<int> handle;
};
\end{lstlisting}

\begin{summarybox}
멤버가 이미 자기 자원을 잘 관리한다면, 생성자/소멸자/복사/이동을 굳이 직접 만들지 않는 편이 더 안전하다.
\end{summarybox}

\section{Factory + Virtual Interface}

\begin{lstlisting}[style=cpp,title={생성과 사용 분리}]
class Optimizer {
public:
    virtual ~Optimizer() = default;
    virtual void step() = 0;
};

std::unique_ptr<Optimizer> createOptimizer(const Config &config) {
    if (config.type == "adam")
        return std::make_unique<Adam>(config.lr);
    return std::make_unique<SGD>(config.lr);
}
\end{lstlisting}

\section{Optional 반환}

\begin{lstlisting}[style=cpp,title={실패 가능성을 타입으로 드러내기}]
std::optional<User> findUser(const std::string &name) {
    auto it = db.find(name);
    if (it != db.end())
        return it->second;
    return std::nullopt;
}
\end{lstlisting}

\section{string\_view / StringRef}

\begin{lstlisting}[style=cpp,title={복사 없이 문자열 빌려보기}]
void log(std::string_view msg) {
    std::cout << msg << "\n";
}
\end{lstlisting}

\begin{warningbox}
\code{string\_view}는 소유하지 않는다. 원본 문자열보다 오래 살면 dangling view가 된다.
\end{warningbox}

\section{constexpr + static\_assert}

\begin{lstlisting}[style=cpp,title={컴파일 타임 계산과 검증}]
constexpr int factorial(int n) {
    return (n <= 1) ? 1 : n * factorial(n - 1);
}

struct PacketHeader {
    uint32_t size;
    uint16_t type;
};

static_assert(std::is_trivially_copyable_v<PacketHeader>,
              "PacketHeader must be trivially copyable");
\end{lstlisting}

\begin{conceptbox}[title={\code{std::is_same\_v}가 뜻하는 것}]
\code{std::is_same\_v<A, B>}는 ``A와 B가 완전히 같은 타입인가''를 나타내는 컴파일 타임 bool 값이다. 함수 호출이 아니라 타입에 대한 질문이다.
\end{conceptbox}

\section{Scope Guard와 Custom Deleter}

\begin{lstlisting}[style=cpp,title={Scope Guard의 최소 형태}]
struct Resource {
    int handle;
};

Resource *acquire();
void release(Resource *r);

struct ScopeGuard {
    std::function<void()> fn;
    ~ScopeGuard() { fn(); }
};

void process() {
    Resource *resource = acquire();
    ScopeGuard guard{[&]() { release(resource); }};
    // 중간에 return하거나 예외가 나도 release 실행
}
\end{lstlisting}

\begin{lstlisting}[style=cpp,title={unique\_ptr + custom deleter}]
auto resource = std::unique_ptr<Resource, decltype(&release)>(
    acquire(), &release
);
\end{lstlisting}

\begin{summarybox}
\begin{itemize}[nosep]
\item cleanup을 일으키는 것은 함수 포인터 자체가 아니라 \term{관리 객체의 소멸자}다.
\item \code{decltype(\&release)}는 \code{release} 함수 포인터의 타입을 자동으로 가져온다.
\item \code{unique_ptr}는 ``자원 + 해제 방법''을 같이 들고 있는 RAII 객체다.
\end{itemize}
\end{summarybox}

\section{CRTP}

\begin{lstlisting}[style=cpp,title={CRTP의 기본 모양}]
template <typename Derived>
class Printable {
public:
    void print() const {
        const auto &self = static_cast<const Derived &>(*this);
        std::cout << self.toString() << "\n";
    }
};

class Dog : public Printable<Dog> {
public:
    std::string toString() const { return "Dog"; }
};
\end{lstlisting}

\begin{conceptbox}[title={왜 쓰는가}]
CRTP는 부모가 자식 타입을 템플릿 인자로 받아서, virtual 없이 자식 메서드에 접근하는 패턴이다. MLIR의 \code{PassWrapper<MyPass, OperationPass<ModuleOp>>}가 대표적인 예다.
\end{conceptbox}

\section{반복자 무효화 방지}

\begin{lstlisting}[style=cpp,title={erase가 다음 반복자를 반환한다}]
for (auto it = vec.begin(); it != vec.end(); ) {
    if (bad(*it))
        it = vec.erase(it);
    else
        ++it;
}
\end{lstlisting}

\section{한 번에 외울 패턴 묶음}

\begin{longtable}{p{0.26\textwidth}p{0.26\textwidth}p{0.38\textwidth}}
\toprule
상황 & 먼저 떠올릴 것 & 대표 패턴 \\
\midrule
\endhead
자원 정리 & 소멸자 & RAII, Scope Guard, unique\_ptr \\
런타임 다형성 & 부모 인터페이스 & virtual + factory \\
컴파일 타임 분기 & 타입 정보 & template, constexpr, type traits \\
복사 줄이기 & view/참조 & const ref, string\_view, StringRef \\
공통 로직 재사용 & 정적 다형성 & CRTP \\
\bottomrule
\end{longtable}

\part{MLIR을 위한 C++}

\chapter{MLIR 큰 그림}

MLIR은 ``IR을 하나만 두는 시스템''이 아니라 \term{여러 dialect와 여러 단계의 lowering을 조합하는 인프라}다.

\section{핵심 개념}

\begin{longtable}{p{0.2\textwidth}p{0.7\textwidth}}
\toprule
개념 & 의미 \\
\midrule
\endhead
Dialect & 연산과 타입의 집합. 예: \code{func}, \code{arith}, \code{linalg} \\
Operation & IR의 기본 단위. 결과 값, 속성, 피연산자, region을 가질 수 있다 \\
Pass & IR을 읽고 변형하는 단계 \\
Pattern Rewriting & 특정 op를 다른 op들로 치환하는 규칙 기반 변환 \\
TableGen & op, type, attr, pass 선언을 기술하는 메타 언어 \\
\bottomrule
\end{longtable}

\begin{conceptbox}[title={왜 MLIR이 C++과 잘 맞는가}]
MLIR은 타입이 풍부하고, 템플릿과 CRTP를 많이 활용하며, 수많은 작은 객체의 수명과 성능을 정교하게 다룬다. 그래서 C++ 기본기가 곧 MLIR 생산성으로 이어진다.
\end{conceptbox}

\section{대표 작업 흐름}

\begin{lstlisting}[style=shell,numbers=none,title={MLIR 개발의 반복 루프}]
TableGen 수정
-> generated headers / cpp 재생성
-> C++ dialect / pass 구현 수정
-> 빌드
-> mlir-opt 또는 내 툴로 테스트
\end{lstlisting}

\chapter{TableGen과 코드 생성}

\section{TableGen은 왜 필요한가}

MLIR 코드는 반복이 많다. op 이름, operands, results, traits, builders, assembly format 같은 정보는 선언적으로 적고, 나머지는 생성 코드에 맡기는 편이 낫다.

\begin{lstlisting}[style=cpp,title={개념적 예시: TableGen이 표현하는 것}]
def My_AddOp : My_Op<"add", [Pure]> {
  let arguments = (ins I32:$lhs, I32:$rhs);
  let results = (outs I32:$result);
  let assemblyFormat = "$lhs `,` $rhs attr-dict `:` type($result)";
}
\end{lstlisting}

\begin{summarybox}
\begin{itemize}[nosep]
\item TableGen은 사람이 반복 선언을 쓰고
\item 생성 코드는 builder, accessor, parser/printer 보일러플레이트를 담당한다
\item hand-written C++는 검증 로직, 리라이트 로직, 특수한 동작에 집중한다
\end{itemize}
\end{summarybox}

\section{실전 감각}

\begin{itemize}[nosep]
\item op 정의를 바꾸면 generated 파일이 함께 바뀐다.
\item 에러가 generated 파일에서 보이더라도 원인은 원래의 \code{.td} 선언인 경우가 많다.
\item 빌드 에러를 읽을 때는 ``내가 선언한 op 시그니처가 생성 코드와 맞는가''를 먼저 본다.
\end{itemize}

\chapter{MLIR 프로젝트의 CMake 구조}

\section{최소 골격}

\begin{lstlisting}[style=cmake,title={MLIR 도구 프로젝트의 최소 뼈대}]
cmake_minimum_required(VERSION 3.20)
project(MyMlirProject LANGUAGES CXX)

find_package(LLVM REQUIRED CONFIG)
find_package(MLIR REQUIRED CONFIG)

list(APPEND CMAKE_MODULE_PATH "${LLVM_CMAKE_DIR}")
list(APPEND CMAKE_MODULE_PATH "${MLIR_CMAKE_DIR}")
include(AddLLVM)
include(AddMLIR)

add_subdirectory(lib)
add_subdirectory(tools)
\end{lstlisting}

\begin{lstlisting}[style=cmake,title={도구 타겟 예시}]
add_executable(my-opt my-opt.cpp)
target_link_libraries(my-opt
  PRIVATE
    MyDialect
    MyConversion
    MLIRIR
    MLIRParser
    MLIRPass
)
\end{lstlisting}

\section{패스 개발 흐름}

\begin{lstlisting}[style=cpp,title={Lowering 패턴의 전형적인 모양}]
class MyOpLowering : public OpConversionPattern<MyOp> {
public:
    using OpConversionPattern<MyOp>::OpConversionPattern;

    LogicalResult matchAndRewrite(
        MyOp op,
        OpAdaptor adaptor,
        ConversionPatternRewriter &rewriter) const override {
        rewriter.replaceOpWithNewOp<arith::AddIOp>(
            op, adaptor.getLhs(), adaptor.getRhs());
        return success();
    }
};
\end{lstlisting}

\begin{lstlisting}[style=cpp,title={PassWrapper를 이용한 패스 껍데기}]
struct LowerMyToArithPass
    : public PassWrapper<LowerMyToArithPass, OperationPass<ModuleOp>> {
    void runOnOperation() override {
        RewritePatternSet patterns(&getContext());
        patterns.add<MyOpLowering>(&getContext());
        // conversion setup ...
    }
};
\end{lstlisting}

\begin{conceptbox}[title={여기서 C++ 개념이 어떻게 연결되는가}]
\begin{itemize}[nosep]
\item \code{OpConversionPattern<MyOp>}는 template다.
\item \code{PassWrapper<Derived, Base>}는 CRTP다.
\item \code{runOnOperation()}은 virtual override다.
\item \code{RewritePatternSet}과 각종 IR 객체는 수명 관리가 중요하므로 참조와 포인터 감각이 필요하다.
\end{itemize}
\end{conceptbox}

\section{MLIR 학습 체크리스트}

\begin{checklistbox}
\begin{itemize}[nosep]
\item dialect, operation, pass, pattern의 역할을 구분할 수 있는가
\item \code{mlir-opt}로 패스를 실행해 볼 수 있는가
\item TableGen 선언과 generated 코드의 관계를 설명할 수 있는가
\item CMake에서 MLIR 라이브러리를 링크하는 최소 구성을 쓸 수 있는가
\item 패턴 하나와 패스 하나를 새로 추가하는 흐름을 머릿속에서 그릴 수 있는가
\end{itemize}
\end{checklistbox}

\chapter{빠르게 익숙해지기 위한 최종 암기표}

\begin{summarybox}
\begin{itemize}[nosep]
\item CMake는 타겟 그래프를 선언하는 도구다.
\item C++를 읽을 때는 문법보다 수명, 소유권, 디스패치 방식을 먼저 본다.
\item 실전에서 가장 자주 보는 패턴은 RAII, unique\_ptr, const ref, optional, factory, virtual, CRTP다.
\item MLIR 코드는 template, CRTP, virtual, generated code가 한 코드베이스에서 동시에 섞여 나온다.
\item 막힐 때는 ``이 객체는 누가 소유하는가'', ``이 호출은 런타임인가 컴파일 타임인가'', ``이 코드는 generated인가 hand-written인가''를 먼저 물어본다.
\end{itemize}
\end{summarybox}

\begin{quizbox}
\textbf{Q1.} \code{unique_ptr<Resource, decltype(\&release)>}에서 cleanup을 실제로 일으키는 것은 무엇인가? \\
\textbf{A.} 함수 포인터의 소멸이 아니라 \code{unique_ptr} 소멸자이다.

\textbf{Q2.} \code{PassWrapper<MyPass, OperationPass<ModuleOp>>}는 왜 CRTP인가? \\
\textbf{A.} 부모 템플릿이 자식 타입 \code{MyPass}를 인자로 받아 compile-time에 자식 타입을 알기 때문이다.

\textbf{Q3.} MLIR CMake에서 가장 먼저 확인할 외부 의존성 선언은 무엇인가? \\
\textbf{A.} 보통 \code{find_package(LLVM REQUIRED CONFIG)}와 \code{find_package(MLIR REQUIRED CONFIG)}다.
\end{quizbox}

\backmatter

\chapter{다음 단계}

이 책을 한 번 읽고 끝내지 말고, 실제 코드에서 반복해서 대응시키는 것이 중요하다.

\begin{enumerate}[nosep]
\item \code{mlir-opt}나 프로젝트 전용 도구의 \code{main()}을 읽어 본다.
\item dialect 하나의 TableGen 정의와 생성 결과를 연결해서 본다.
\item lowering 패턴 하나를 처음부터 끝까지 추적한다.
\item CMake에서 타겟 하나가 어떤 라이브러리를 타고 들어가는지 따라가 본다.
\end{enumerate}

이 과정을 세 번 정도 반복하면, C++ 문법이 아니라 코드 구조가 먼저 보이기 시작한다.

\end{document}
